\chapter{ПРОГРАМНЕ ЗАБЕЗПЕЧЕННЯ}

%Hidden cite for correct bibliography ordering
\nocite{bahvalov-et-al,benerdge-et-al} 

\section{Огляд проблеми}

Для отримання, валідації і оформлення наукових результатів було створене середовище із наступними властивостями:
\begin{itemize}
    \item Легке встановлення на кількох різних машинах.
    \item Контрольовані зміни до налаштувань, які легко поширити на всі машини.
    \item Зручна візуалізація даних.
    \item Можливість працювати із абстракціями близькими до математичного апарату.
\end{itemize}

\noindent Стабільне середовище, яке дозволяє проводити дослідження як на потужних так і на мобільних машинах, дозволяє 
використати час більш продуктивно. З іншої сторони, наявність інтерфейсу високого рівня, який дозволяє працювати із 
формулами майже у співвідношенні один до одного зменшує час на пошуки помилок і відлагодження програмного забезпечення. 

\section{Середовище}

Технологія Docker дозволяє описувати середовище виконання у вигляді коду. Це середовище, яке називається Docker 
відображенням, можна завантажити на спеціальний сервер. На довільній машині із операційною системою Ubuntu можна 
вивантажити це відображення і запустити відповідний контейнер. Він буде містити ідентичне середовище, яке буде 
виконуватися ідентично (при умові, що версії ядер Ubuntu однакові). Також Docker відображенням можна збудувати на 
довільній машині використовуючи текстовий файл опису Dockefile. Це суттєва перевага цієї технології над іншими. Завдяки
таким властивостям можна завантажити Dockerfile у систему зберігання коду і відслідковувати зміни. Це дозволяє будувати
Docker відображення на локальній машині із Dockefile, у випадку, коли немає сервера із потрібним Docker відображенням.

На (рис. \ref{fig:software_docker_hierarchy}) зображена побудована ієрархія Docker відображень. Вона доволі гнучка, бо
дозволяє використовувати відображення із потрібним набором програмного забезпечення на конкретній машині.

\begin{figure}[ht!]
    \centering
    \begin{tikzpicture}
        \begin{class}[text width=2cm]{BaseUI}{0,0}
        \end{class}

        \begin{class}[text width=2cm]{NGSolve}{0,-2}
            \inherit{BaseUI}
        \end{class}

        \begin{class}[text width=2cm]{NGBem}{0,-4}
            \inherit{NGSolve}
        \end{class}

        \begin{class}[text width=2cm]{LaTeX}{0,-6}
            \inherit{NGBem}
        \end{class}

    \end{tikzpicture}
    \caption{Ієрархія Docker відображень}
    \label{fig:software_docker_hierarchy}
\end{figure}


\section{Програмні продукти}

Для виконання обчислень використовувалася програмний пакет NGSolve \cite{software-ngsolve} і його розширення NGBem.

\section{Висновок}

Засоби дозволяють все зберігати у текстових файлах.


Протягом багатьох років підходи до програмування систем мінялися в залежності від потреб часу. Спочатку програми
складалися на машинній мові, пізніше виникли асемблери, які давали змогу транслювати команди зрозумілі людині в машинні
коди. Щоб спростити написання програми і наблизити мови програмування до людської мови виникли макропроцесори, які
перетворюють макрокоманди у команди асемблера. Далі із мікроасемблерних бібліотек з'явилися мови процедурного типу, які
мали власний транслятор у машинні коди.

Процедурні мови програмування набули широкого розповсюдження і використовувалися для побудови обчислювальних алгоритмів.
Найбільшими проблемами, з якими стикаються люди у цих мовах програмування є невисока аналогія із реальними об'єктами,
які моделюються. Процедури дозволяються реалізувати алгоритми і розбити їх на підзадачі.

Об'єктно-орієнтоване програмування - одна з парадигм програмування, на відміну від процедурного програмування дозволяє
ширші можливості застосування. У процедурному ("алгоритмічному") програмуванні (програмування за допомогою мов Fortran,
Бейсік, Pascal та їм подібних "алгоритмічних" мов) уся увага зосереджується на розробці та проробці системи взаємодіючих
процедур та функцій, часто згрупованих у модулі та бібліотеки за семантичними та іншими ознаками, котрі реалізують
алгоритми необхідні для функціонування програми або операційної системи. Дані у таких програмах зберігаються у
глобальних (відносно окремих процедур) змінних і передаються у процедури та функції як параметри.

На відміну від того в об'єктно-орієнтованому програмуванні (використовують мови Simula-67, С++, Smalltalk, Java, Python
та їм подібні) дані та методи (процедури) пов'язані з ними є комбінованими у класи (об'єкти), що відповідають
онтологічним сутностям прикладної області або є допоміжними у програмі. Тоді конкретні значення полів(змінних) об'єкта
визначають його стан, а його методи дозволяють іншим об'єктам програми взаємодіяти з ним. Доступні методи об'єкта ще
називають його інтерфейсом (контрактом).

ООП базується на трьох парадигмах: інкапсуляція, успадкування, поліморфізм.

\begin{enumerate}
\def\labelenumi{\arabic{enumi}.}
\tightlist
\item
  Інкапсуляція - об'єкт приховує свою внутрішню реалізацію, залишаючи
  для користувачів інтерфейси взаємодії
\item
  Успадкування - нащадок об'єкта отримує ті ж самі властивості, що й
  об'єкт-предок.
\item
  Поліморфізм - об'єкт змінює свою поведінку у рамках інтерфейсу,
  успадкованому від батьківських об'єктів
\end{enumerate}

Основні проблеми програмного забезпечення С++

\begin{enumerate}
\def\labelenumi{\arabic{enumi}.}
\tightlist
\item
  Вказівники
\item
  Бібліотечні класи перевантажені функціональністю
\item
  При процедурному стилі внутрішні дані видимі всім
\item
  Різні компілятори мають свої особливості
\end{enumerate}

Переваги

\begin{enumerate}
\def\labelenumi{\arabic{enumi}.}
\tightlist
\item
  С++ мова, компілятори для якої є на багатьох платформах
\item
  Багатошарова структура (розбиття на модулі) дозволяє рознесення різних
  частин програми на різні фізичні машини для підвищення продуктивності
  системи
\item
  Внутрішні дані потрібні для певного об'єкту залишаються недосяжними
  для інших об'єктів
\item
  Реалізовано ряд алгоритмів, які вже не треба повторно писати
\item
  Масштабованість -- без зміни існуючого коду можна додати новий код
\item
  Багатошарова структура програми дозволяє міняти один із шарів без
  зміни іншого
\item
  Дані, що виводяться можна без особливих зусиль писати у довільному
  форматі, оскільки сам розв'язок задачі подається у зручному вигляді
\end{enumerate}

Патерни застосовуються для того щоб описати найтиповіші рішення.

Система складається із наступних підсистем, які можна зобразити у
вигляді модулів:

\begin{enumerate}
\def\labelenumi{\arabic{enumi}.}
\tightlist
\item
  Введення даних
\item
  Класів загального використання
\item
  Обчислень
\item
  Виведення даних
\end{enumerate}
